\chapter{Introduction}
\label{ch:Introduction}

The description of a software system is typically spread across several models of different languages \cite{KLARE2021110815}.
A UML class diagram typically represents the system's design, while the Java source code contains its implementation.
Both artifacts describe the same classes, resulting in overlapping information.
Renaming a class in the UML model necessitates renaming the corresponding Java class and its associated file.
When such a change is not carried over, the models contradict each other, and the system description becomes inconsistent \cite{KLARE2021110815}.
Manual consistency checking and restoration are error-prone, and resolving inconsistencies can be costly \cite{KLARE2021110815}.

\section{Motivation}
\label{sec:Introduction:Motivation}

Consistency preservation rules automate this process by propagating changes made in one model to all related models.
Vitruvius is a framework that extends these rules \cite{KLARE2021110815}.
Vitruvius integrates system models into a Virtual Single Underlying Model (V-SUM), which developers access and modify through specific views.
Following each modification, the V-SUM applies consistency preservation rules specified by a methodologist using a dedicated language, such as the Reactions Language \cite{KLARE2021110815}.
Developers interact with the models but are not permitted to alter the consistency preservation rules \cite{KLARE2021110815}.
The methodologist determines which elements the V-SUM derives and their respective forms during the rule specification process.
Vitruvius and the Reactions Language are discussed in detail in \autoref{sec:Foundations:Vitruvius} and \autoref{sec:Foundations:ReactionsLanguage}.

A consistency specification commits to a single mapping between the involved metamodels, even though multiple mappings may be justifiable.
A UML design can be mapped to Java in various ways to accommodate different implementation constraints \cite{harrison2000mapping}.
For example, Harrison et al. represent each UML class as a Java interface with two implementing classes, exposing each attribute exclusively through getter and setter methods \cite{harrison2000mapping}.
Certain decisions are determined by platform specific requirements.
A JavaBeans component defines a property by supplying public getter and setter methods \cite{javaBeansTutorial}.
A Jakarta Persistence entity must be implemented as a class, not as an interface or enumeration \cite{jakartaPersistence}.
That class must not be final and needs a public or protected constructor without parameters \cite{jakartaPersistence}.
The appropriate realization produced by a rule set therefore depends on the specific requirements of the project in which it is applied.

A methodologist who wants to support several conventions currently has to maintain a separate copy of the rule set for each of them.
A correction to a rule that the copies share then has to be applied to every copy, and a copy that misses it diverges from the others.
Since every combination of independent choices requires a rule set of its own, the number of copies multiplies with each further choice.

Furthermore, the rule set applied to a project is subject to change over time.
Conventions may evolve throughout the system's lifecycle, such as when a project transitions to a new Java version or adopts updated guidelines.
Rules are also revised in response to identified defects, as demonstrated by the corrections made to the original umljava rules during this work (\autoref{sec:Evaluation:AnnotationMechanism:ChangesToTheOriginalRules}).
At that point, V-SUMs constructed under the previous rule set are already present.
A persisted V-SUM holds its models, the correspondences between their elements and the identifiers of those elements, but no information about the rules that produced them \cite{vitruviusGithub}.
When loaded with a new rule set, the V-SUM preserves all elements derived by the previous rules, while only subsequent modifications adhere to the updated rules.
This process results in a V-SUM that is not consistent with either the original or the updated rule set.
In the absence of a migration mechanism, a project must either manually repair its V-SUM, regenerate it entirely under the new rules, or continue using the previous rules.
A derivation from scratch discards the handwritten content that no rule derives.
Klare et al. identify the evolution of a V-SUM metamodel, including its consistency preservation rules, as an area for future research \cite{KLARE2021110815}.
Each V-SUM constructed with such a metamodel must then co-evolve to maintain consistency among its models \cite{KLARE2021110815}.
No publication was found that migrates a persisted V-SUM to a changed consistency specification (\autoref{sec:RelatedWork:ResearchGap}).

This thesis addresses both problems.
In accordance with the principles of software product lines, each reaction is annotated by a methodologist with its corresponding feature, ensuring that a single rule set encompasses the reactions of all conventions (\autoref{sec:Foundations:SoftwareProductLines}).
This rule set is called the 150\% rule set.
Project-specific rules are derived from the 150\% rule set based on a given configuration, which represents the set of features selected for the project.
A derived rule set contains no annotations, so the compiler of the Reactions Language processes it like any ordinary rule set.
When the rules of an existing V-SUM change, a migration brings the V-SUM to the new rules.
The migration process either retains a single dominant model and re-derives all other models from it, or limits modifications to only those elements affected by the added, removed, or changed rules.
Content without correspondence, such as handwritten method bodies, can be identified and subsequently restored.

\section{Research goals}
\label{sec:Introduction:ResearchGoals}

This thesis addresses two primary objectives.
The first objective is to make the consistency preservation rules of a V-SUM variable by annotating reactions with features.
The second objective is to migrate a persisted V-SUM to a changed rule set, so that it becomes consistent under the new rules.
Both objectives are evaluated using the Goal Question Metric (GQM) approach, which refines each goal into quantifiable questions and addresses each question with a corresponding set of metrics \cite{basili1992gqm}.
G1 focuses on the annotation process and the derivation of the rule set.
G2 through G4 address the migration process, specifically evaluating the quality of the migration outcome, identifying the impact of rule changes, and assessing the required effort.
G5 investigates the applicability of the approach to rule sets beyond those for which it was originally developed.
Most questions are answered on a library example that is derived with the rules of the umljava case study \cite{vitruv-casestudies-fork} in each of their nine configurations (\autoref{ch:RunningExample}).
The four metamodels of BrakeCaseStudy are additionally used to compare strategies for selecting the dominant model.
For G5, the annotation mechanism is applied to the pcmumlclass rules.

\begin{description}
      \item[G1: Annotation mechanism.] \label{gqm:G1} Analyze the annotation mechanism with respect to compactness, adaptation effort, the changes to the original rules, validity of the derived configurations, and processing overhead.
            \begin{itemize}
                  \item[Q1.1] How compact is a 150\% rule set compared to maintaining the rule set of every configuration separately?
                  \item[M1.1] LOC and number of files of the 150\% rule set compared to the sum over all separately maintained configurations.
                  \item[Q1.2] How much adaptation effort is required when a feature or a configuration is added or changed?
                  \item[M1.2] Number of files created or modified per added or changed feature or configuration.
                  \item[Q1.3] How much did the original rule set have to be changed to become a 150\% rule set?
                  \item[M1.3] Lines of code the 150\% rule set adds to the original rule set, split into annotations, new reactions and routines, and changes to existing reactions and routines, the latter classified by cause.
                  \item[Q1.4] Do the configurations derived by the preprocessor produce valid and operational rules?
                  \item[M1.4] Per derived configuration whether the generated Java compiles and how many case study tests pass with and without their feature gates, and the number of original tests passing on config1.
                  \item[Q1.5] How much overhead does preprocessing add when a configuration is built and tested?
                  \item[M1.5] Preprocessing time, in absolute terms and as a share of a test run of the case study.
            \end{itemize}
      \item[G2: Migration quality.] \label{gqm:G2} Analyze the migration with respect to its equivalence to a derivation from scratch, the validity of its artifacts, the content it carries over, its reversibility, and the choice of the dominant model.
            \begin{itemize}
                  \item[Q2.1] Does a migrated V-SUM describe the same state that the new rules derive from scratch?
                  \item[M2.1] Number of configuration pairs whose migrated models are equivalent to a derivation of the target configuration from scratch.
                  \item[Q2.2] Are the artifacts produced by a migration valid?
                  \item[M2.2] Number of migration runs whose resulting Java compiles.
                  \item[Q2.3] How much content cannot be carried over automatically?
                  \item[M2.3] Number of handwritten elements the preservation report lists as preservable or lost per migration, and the number of the other listed elements grouped by cause.
                  \item[Q2.4] Is a migration reversible?
                  \item[M2.4] Number of configuration pairs for which migrating back reproduces the original V-SUM.
                  \item[Q2.5] How do the strategies for choosing the dominant model compare?
                  \item[M2.5] Chosen dominant model, resulting number of changes, and selection time per strategy.
            \end{itemize}
      \item[G3: Inconsistency identification.] \label{gqm:G3} Analyze the comparison of rule sets and the matching of triggers with respect to the inconsistencies a rule change causes, the precision with which affected elements are identified, and the rule changes only semantic hashes detect.
            \begin{itemize}
                  \item[Q3.1] Which inconsistencies can arise when a feature selection is changed after artifacts already exist?
                  \item[M3.1] Number and type of identified inconsistency classes.
                  \item[Q3.2] How precisely are the affected model elements identified?
                  \item[M3.2] Reported dirty rules and affected elements compared to the elements actually differing from a derivation from scratch.
                  \item[Q3.3] Which rule changes require a semantic hash instead of a comparison of reaction identifiers?
                  \item[M3.3] Number of kinds of rule changes detected by comparing reaction identifiers and by comparing semantic hashes.
            \end{itemize}
      \item[G4: Migration effort.] \label{gqm:G4} Analyze the selective migration with respect to its scope and runtime compared to a full migration.
            \begin{itemize}
                  \item[Q4.1] Which share of the rules and model elements of a V-SUM does a selective migration process?
                  \item[M4.1] Dirty rules relative to all rules and repropagated elements relative to all elements of the dominant model, per configuration pair.
                  \item[Q4.2] Does the reduced scope translate into a reduced runtime?
                  \item[M4.2] Runtime of the selective compared to the full migration.
            \end{itemize}
      \item[G5: Generalizability.] \label{gqm:G5} Analyze the approach with respect to transferability to rules beyond the ones it was developed on.
            \begin{itemize}
                  \item[Q5.1] Does the transferred solution remain operational?
                  \item[M5.1] Number of existing tests that pass without changes.
                  \item[Q5.2] Is the mechanism independent of the concrete rules and metamodels it is applied to?
                  \item[M5.2] Number of specific code paths in the preprocessor and the migration.
            \end{itemize}
\end{description}

\section{Contributions}
\label{sec:Introduction:Contributions}

This thesis makes four contributions.

\begin{itemize}
      \item Feature annotations in the Reactions Language associate specific reactions with corresponding features.
            A preprocessor subsequently derives the rule set for a given configuration from a 150\% rule set (\autoref{sec:Concept:FeatureAnnotations}, \autoref{sec:Implementation:Preprocessor}).
            Because an annotation retains the declaration of a reaction, each rule maintains its identifier in every configuration in which it appears.
            This property serves as the foundation for selective migration.
      \item The rules of the umljava case study are expanded into a 150\% rule set, with a feature model comprising 37 features.
            From this model, nine distinct configurations are derived (\autoref{ch:RunningExample}).
            The rule set introduces alternative realizations for UML classes and data types, incorporates optional conventions, and addresses defects present in the original rules.
            The tests within the case study are gated by the features they require or exclude (\autoref{sec:Implementation:Preprocessor}, \autoref{sec:Implementation:Testing}).
      \item Migration applies an updated rule set to a persisted V-SUM using one of two approaches (\autoref{sec:Concept:Migration}, \autoref{sec:Implementation:Migration}).
            Full migration retains a dominant model, selected through one of three strategies, and subsequently derives all other models from it.
            Selective migration compares the rules recorded in the V-SUM with the updated rules and repropagates only the affected elements, the elements that the added, removed or changed rules apply to.
            To support this process, change propagation in Vitruvius stores the identifier, semantic hash, and a trigger summary for each rule.
            Additionally, a defect in change recording for cascaded deletions has been addressed \cite{vitruv-change-issue}.
            A subsequent preservation step reports or restores underived content, specifically content not referenced by any correspondence, which would otherwise be lost during migration.
      \item The approach is evaluated according to the GQM plan using 81 ordered configuration pairs from the library example, the BrakeCaseStudy, and a transfer to the pcmumlclass rules (\autoref{ch:Evaluation}).
            Of these pairs, 72 switch between two different configurations, whereas the other nine keep their configuration.
            The generated artifacts and evaluation scripts have been made publicly available \cite{thesis-repo}.
\end{itemize}

The implementation extends the Vitruvius repositories Vitruv-DSLs \cite{vitruv-dsls-fork}, Vitruv-Change \cite{vitruv-change-fork} and Vitruv-CaseStudies \cite{vitruv-casestudies-fork}.
In addition, BrakeCaseStudy is ported to Vitruvius 4 together with the extended Vitruv-Change and Vitruv-DSLs \cite{brakecasestudy-fork}.

\section{Thesis structure}
\label{sec:Introduction:ThesisStructure}

This thesis begins by establishing foundational concepts, reviewing related work, and introducing the running example.
It then presents the proposed approach and its implementation, followed by evaluation, conclusion, and directions for future research.
The foundational section in \autoref{ch:Foundations} addresses the Eclipse Modeling Framework, Vitruvius, the Reactions Language, JaMoPP, and software product lines.
\Autoref{ch:RelatedWork} discusses related approaches to variability, migration of existing artifacts, preservation of handwritten content, and consistency.
The chapter concludes by identifying the research gap.
The running example in \autoref{ch:RunningExample} presents the umljava rules, their associated feature model, and nine configurations, as well as the library example used in the evaluation.
The conceptual framework in \autoref{ch:Concept} details the feature annotations, derivation of the rule set, procedures for both full and selective migration, and the preservation process.
The implementation described in \autoref{ch:Implementation} includes the preprocessor, migration tool, testing procedures, and the matrix run that generates the majority of evaluated data.
The evaluation in \autoref{ch:Evaluation} addresses the research questions formulated in the Goal-Question-Metric (GQM) plan.
Finally, \autoref{ch:Conclusion} summarizes the main findings, while \autoref{ch:FutureWork} proposes directions for future research.
